\chapter{Machine Learning Foundations \& Data Preprocessing}

\section{Challenges in Machine Learning}

Machine learning applications in financial engineering face several structural and regulatory constraints:

\begin{itemize}[leftmargin=1.5em]
    \item \textbf{Requires Large, Clean Data}: Models demand substantial volumes of high-quality, cleansed historical records to generalize effectively.
    \item \textbf{Risk of Overfitting}: Complex models risk fitting noise and idiosyncratic market movements rather than genuine underlying signals.
    \item \textbf{Interpretability (\textquotedblleft Black Box\textquotedblright)}: Advanced algorithms often lack transparent mechanistic explanations, creating auditing hurdles.
    \item \textbf{Regulatory \& Ethical Concerns in Finance}: Strict regulatory compliance mandates model explainability, fairness, and accountability in financial decision-making.
\end{itemize}

\section{Types of Variables}

Variables encountered in machine learning datasets are broadly categorized into numerical and categorical types:

\begin{center}
\begin{tikzpicture}[
    level 1/.style={sibling distance=7.5cm, level distance=1.3cm},
    level 2/.style={sibling distance=3.2cm, level distance=1.3cm},
    every node/.style={
        draw=navyblue,
        thick,
        rounded corners=2pt,
        fill=lightgraybg,
        font=\sffamily\bfseries\small,
        text=navyblue,
        inner sep=5pt,
        align=center
    },
    edge from parent/.style={draw=navyblue, -{Stealth[length=2.5mm]}, thick}
]
\node {Variables}
    child { node {Numerical}
        child { node {Continuous} }
        child { node {Discrete} }
    }
    child { node {Categorical}
        child { node {Nominal} }
        child { node {Ordinal} }
        child { node {Binary} }
    };
\end{tikzpicture}
\end{center}

\begin{table}[htbp]
\centering
\caption{Taxonomy of Data Types and Financial Examples}
\vspace{0.2em}
\begin{tabularx}{\textwidth}{llp{4.2cm}X}
\toprule
\textbf{Type} & \textbf{Subtype} & \textbf{Meaning} & \textbf{Finance Example} \\
\midrule
Numerical & Continuous & Any value in a continuous range & Stock price, portfolio return, volatility \\
\addlinespace[0.4em]
Numerical & Discrete & Countable whole numbers & Number of trades per day, number of defaults \\
\addlinespace[0.4em]
Categorical & Nominal & No natural ordering among categories & Sector: Banking, IT, Pharma \\
\addlinespace[0.4em]
Categorical & Ordinal & Natural hierarchical order exists & Credit Rating: $\text{AAA} > \text{AA} > \text{A} > \text{BBB}$ \\
\addlinespace[0.4em]
Categorical & Binary & Exactly two distinct categories & Loan Default: Yes / No \\
\bottomrule
\end{tabularx}
\end{table}

\section{Data Preprocessing Pipeline}

The standard end-to-end data preparation workflow transforms raw financial inputs into model-ready datasets:

\begin{figure}[htbp]
\centering
\begin{tikzpicture}[
    stage/.style={
        draw=navyblue,
        line width=1.1pt,
        fill=navyblue!5,
        font=\sffamily\bfseries\small,
        text=navyblue,
        rounded corners=3pt,
        text width=2.4cm,
        minimum height=1.1cm,
        align=center
    },
    finalstage/.style={
        draw=darkteal,
        line width=1.3pt,
        fill=darkteal!8,
        font=\sffamily\bfseries\small,
        text=darkteal!90!black,
        rounded corners=3pt,
        text width=2.8cm,
        minimum height=1.1cm,
        align=center
    },
    conn/.style={
        draw=navyblue,
        -{Stealth[length=2.8mm, width=2.0mm]},
        line width=1.2pt
    }
]
\node[stage] (raw) at (0, 0) {Raw Data};
\node[stage] (clean) at (3.3, 0) {Cleaning};
\node[stage] (trans) at (6.6, 0) {Transformation};
\node[stage] (norm) at (9.9, 0) {Normalization /\\Scaling};
\node[stage] (enc) at (9.9, -2.0) {Encoding\\Categorical};
\node[finalstage] (ready) at (6.6, -2.0) {Model-Ready\\Data};

\draw[conn] (raw) -- (clean);
\draw[conn] (clean) -- (trans);
\draw[conn] (trans) -- (norm);
\draw[conn] (norm) -- (enc);
\draw[conn] (enc) -- (ready);
\end{tikzpicture}
\caption{End-to-End Data Preprocessing Pipeline from Raw Inputs to Model-Ready Data.}
\label{fig:preprocessing_pipeline}
\end{figure}

\begin{itemize}[leftmargin=1.5em]
    \item \textbf{Cleaning}: Removes duplicate records, corrects erroneous entries, and handles missing observations.
    \item \textbf{Transformation}: Converts raw figures into mathematically suitable modeling formats (e.g., calculating log returns from sequential stock prices).
    \item \textbf{Normalization / Scaling}: Brings numerical features across heterogeneous magnitudes into a common comparable scale.
    \item \textbf{Encoding}: Converts non-numeric categorical attributes into numerical indicator representations.
\end{itemize}

\subsection{Scaling \& Normalization}

Features measured in disparate units must be rescaled to ensure distance-based and gradient-based algorithms treat features proportionately:

\begin{formulabox}[Min-Max Scaling]
Rescales feature values into the bounded closed interval $[0, 1]$:
\begin{equation}
x' = \frac{x - x_{\min}}{x_{\max} - x_{\min}}
\end{equation}
\end{formulabox}

\begin{formulabox}[$Z$-Score Standardization]
Standardizes observations using empirical sample mean and sample standard deviation:
\begin{equation}
Z = \frac{x - \mu}{\sigma}
\end{equation}
\end{formulabox}

\section{Train-Test-Validation Splitting}

Datasets are partitioned into distinct subsets to evaluate model performance on unseen data, verifying whether the algorithm generalizes rather than memorizing training samples:

\begin{itemize}[leftmargin=1.5em]
    \item \textbf{Training Set}: Used to fit model parameters and learn patterns.
    \item \textbf{Validation Set}: Used during model development to tune hyperparameters and guide model selection.
    \item \textbf{Testing Set}: Used strictly for final unbiased model performance evaluation.
\end{itemize}

\subsection{Common Partitioning Ratios}
\begin{itemize}[leftmargin=1.5em]
    \item \textbf{Train-Test Splits}: $80\% - 20\%$ or $70\% - 30\%$.
    \item \textbf{Train-Validation-Test Splits}: $70\% - 15\% - 15\%$.
\end{itemize}

\subsection{K-Fold Cross-Validation}
In $k$-fold cross-validation:
\begin{enumerate}[leftmargin=1.5em]
    \item The training dataset is partitioned into $k$ equal parts (\textquotedblleft folds\textquotedblright).
    \item The model is trained on $k-1$ folds and validated on the remaining fold.
    \item This procedure is repeated $k$ times, rotating the held-out validation fold sequentially across all splits.
\end{enumerate}

\begin{remarkbox}[Benefits \& Common Choices]
\begin{itemize}[leftmargin=1.2em]
    \item Provides a significantly more robust estimate of true out-of-sample model performance.
    \item Reduces vulnerability to an anomalous single lucky or unlucky split.
    \item Common configurations: $k = 5$ or $k = 10$.
\end{itemize}
\end{remarkbox}

\subsection{Data Leakage Prevention}
\begin{takeawaybox}[Data Leakage Warning]
\textbf{Never scale or normalize the entire dataset prior to splitting.}

If sample mean $\mu$ and standard deviation $\sigma$ are calculated over the aggregate dataset, information from the test set leaks directly into the training stage. 

\textbf{Mandatory Pipeline Sequence:}
\begin{enumerate}[leftmargin=1.2em]
    \item Split the raw dataset first.
    \item Fit the scaler/imputer exclusively on training data.
    \item Transform validation and test data using the saved training parameters.
\end{enumerate}
\end{takeawaybox}

\subsection{Splitting Financial Time-Series Data}
For time-series data like stock prices, a purely random split is fundamentally problematic because it mixes past and future observations. Sequential financial records must strictly use \textbf{chronological splits}.

\section{Missing Data Mechanisms \& Imputation}

\subsection{Missingness Mechanisms}

\begin{table}[htbp]
\centering
\caption{Mechanisms of Missing Data and Practical Financial Examples}
\vspace{0.2em}
\begin{tabularx}{\textwidth}{p{2.0cm}p{6.0cm}X}
\toprule
\textbf{Type} & \textbf{Meaning} & \textbf{Finance Example} \\
\midrule
\textbf{MCAR} & \textbf{Missing Completely at Random}: Missingness is entirely independent of observed and unobserved variables. & Trade entry terminal intermittently experiences random network packet drops logging trades. \\
\addlinespace[0.5em]
\textbf{MAR} & \textbf{Missing at Random}: Missingness is statistically related to other observed features in the dataset. & Younger retail banking customers are systematically more likely to leave the optional income field blank. \\
\addlinespace[0.5em]
\textbf{MNAR} & \textbf{Missing Not at Random}: Missingness is directly correlated with the actual missing value itself. & Highly indebted distressed borrowers deliberately omit outstanding secondary credit obligations. \\
\bottomrule
\end{tabularx}
\end{table}

\subsection{Imputation Methods}

Common techniques for replacing missing data include:
\begin{itemize*}[itemjoin={, }, itemjoin*={, and }]
    \item Mean / Median / Mode
    \item K-Nearest Neighbors (KNN) Imputation
    \item Forward Fill
    \item Backward Fill
    \item Interpolation
\end{itemize*}.

\begin{table}[htbp]
\centering
\caption{Advantages and Disadvantages of Imputation Methods}
\vspace{0.2em}
\begin{tabularx}{\textwidth}{p{3.2cm}p{4.8cm}X}
\toprule
\textbf{Method} & \textbf{Advantages} & \textbf{Disadvantages} \\
\midrule
Mean / Median & Simple, fast to compute & Distorts sample variance; ignores inter-feature relationships. \\
\addlinespace[0.4em]
Mode & Directly applicable to categorical data & Can artificially over-represent the most common category. \\
\addlinespace[0.4em]
KNN & Exploits multivariate relationships between features & Computationally expensive on large financial datasets. \\
\addlinespace[0.4em]
Forward / Backward Fill & Well-suited for time-series sequences & Assumes most recent observed value remains representative. \\
\addlinespace[0.4em]
Interpolation & Produces smooth estimates capturing trends & Unsuitable for non-sequential, unordered data. \\
\bottomrule
\end{tabularx}
\end{table}

\subsection{Decision Framework for Selecting Imputation Methods}

\begin{figure}[htbp]
\centering
\begin{tikzpicture}[
    decision/.style={
        draw=navyblue,
        line width=1.2pt,
        fill=navyblue!6,
        font=\sffamily\bfseries\small,
        text=navyblue,
        rounded corners=3pt,
        align=center,
        inner sep=6pt,
        text width=4.2cm,
        minimum height=1.0cm
    },
    action/.style={
        draw=darkteal,
        line width=1.2pt,
        fill=darkteal!6,
        font=\sffamily\small,
        text=darkteal!90!black,
        rounded corners=3pt,
        align=center,
        inner sep=6pt,
        minimum height=1.0cm
    },
    conn/.style={
        draw=navyblue,
        -{Stealth[length=2.8mm, width=2.0mm]},
        line width=1.2pt
    },
    badge/.style={
        font=\footnotesize\bfseries\sffamily,
        fill=white,
        draw=navyblue!40,
        text=navyblue,
        rounded corners=2pt,
        inner sep=2.5pt
    }
]
% Level 0: Root Question
\node[decision] (ts) at (0, 0) {Is data Time Series?};

% Level 1: Branches
\node[action, text width=4.0cm] (ts_act) at (-4.0, -2.0) {Forward / Backward Fill\\or Interpolation};
\node[decision, text width=4.0cm] (dtype) at (4.0, -2.0) {Numerical or\\Categorical?};

% Level 2: Data Type Leaves
\node[action, text width=3.4cm] (num_act) at (2.0, -4.4) {Mean / Median /\\KNN Imputation};
\node[action, text width=3.0cm] (cat_act) at (6.0, -4.4) {Mode\\Imputation};

% Connectors with ample space
\draw[conn] (ts) -- node[above left=2pt, badge] {Yes} (ts_act);
\draw[conn] (ts) -- node[above right=2pt, badge] {No} (dtype);
\draw[conn] (dtype) -- node[above left=2pt, badge] {Numerical} (num_act);
\draw[conn] (dtype) -- node[above right=2pt, badge] {Categorical} (cat_act);
\end{tikzpicture}
\caption{Decision Framework for Selecting Missing Value Imputation Methods.}
\label{fig:imputation_decision_framework}
\end{figure}

\section{Imbalanced Categorical Data}

Class imbalance arises when one category of the target variable appears far more frequently than other categories:

\begin{remarkbox}[Financial Example]
In credit card fraud detection:
\begin{itemize}[leftmargin=1.5em]
    \item Genuine Transactions: $99.9\%$
    \item Fraudulent Transactions: $0.1\%$
\end{itemize}
A naive model that blindly predicts \textquotedblleft no fraud\textquotedblright\ every single time achieves $99.9\%$ raw accuracy, yet is entirely useless. Thus, \textbf{accuracy alone is fundamentally misleading}.
\end{remarkbox}

\subsection{Consequences of Class Imbalance}
\begin{itemize}[leftmargin=1.5em]
    \item Biases model parameters heavily toward majority class patterns.
    \item Renders overall accuracy metric uninformative for decision-makers.
\end{itemize}

\subsection{Remedial Solutions}
\begin{itemize}[leftmargin=1.5em]
    \item \textbf{Oversampling}: Duplicating or synthetically generating minority class instances.
    \item \textbf{Undersampling}: Down-sampling majority class observations.
    \item \textbf{SMOTE}: Synthetic Minority Over-sampling Technique, interpolating new instances along minority feature lines.
\end{itemize}

\section{Model Evaluation Metrics}

For binary classification problems, predictions and actual outcomes are summarized in a confusion matrix:

\begin{table}[htbp]
\centering
\caption{Binary Classification Confusion Matrix}
\vspace{0.2em}
\begin{tabular}{lcc}
\toprule
& \textbf{Predicted Positive} & \textbf{Predicted Negative} \\
\midrule
\textbf{Actual Positive} & True Positive ($TP$) & False Negative ($FN$) \\
\addlinespace[0.3em]
\textbf{Actual Negative} & False Positive ($FP$) & True Negative ($TN$) \\
\bottomrule
\end{tabular}
\end{table}

\begin{formulabox}[Precision]
Of all instances predicted as positive, how many were genuinely correct?
\begin{equation}
\text{Precision} = \frac{TP}{TP + FP}
\end{equation}
\end{formulabox}

\begin{formulabox}[Recall]
Of all actual positive instances in the population, how many did the model identify?
\begin{equation}
\text{Recall} = \frac{TP}{TP + FN}
\end{equation}
\end{formulabox}

\begin{formulabox}[$F_1$-Score]
Harmonic mean of Precision and Recall, balancing false alarms against missed detections:
\begin{equation}
F_1 = 2 \cdot \frac{\text{Precision} \cdot \text{Recall}}{\text{Precision} + \text{Recall}}
\end{equation}
\end{formulabox}

\section{Mistakes to be Avoided in ML Pipelines}

\begin{enumerate}[leftmargin=1.5em]
    \item \textbf{Global Scaling Before Splitting}: Scaling or normalizing the entire dataset prior to train/test partitioning causes data leakage.
    \item \textbf{Blind Mean Imputation on MNAR Data}: Imputing missing values with unconditional means when data is Missing Not at Random injects substantial estimation bias.
    \item \textbf{Ignoring Class Imbalance}: Relying exclusively on overall classification accuracy in skewed datasets conceals poor minority detection.
    \item \textbf{Unchecked One-Hot Encoding on High-Cardinality Variables}: Creating dummy variables for high-cardinality categorical features leads to severe dimensionality inflation.
    \item \textbf{Failing to Shuffle Non-Time-Series Data}: Neglecting random shuffling before partitioning standard cross-sectional data risks clustering-induced bias.
\end{enumerate}

\section{Key Takeaways}

\begin{takeawaybox}[Core Principles of ML Data Preparation]
\begin{enumerate}[leftmargin=1.2em, label=\roman*)]
    \item \textbf{Pre-processing dominates}: Data preprocessing accounts for $60\% - 80\%$ of total machine learning workflow effort.
    \item \textbf{Always split first}: Partition data into training, validation, and test sets before performing any scaling or modeling.
    \item \textbf{Diagnose missingness}: Determine whether missing values are MCAR, MAR, or MNAR before selecting an imputation strategy.
    \item \textbf{Address class imbalance}: Use SMOTE, oversampling, or undersampling; accuracy alone is never sufficient.
    \item \textbf{Respect temporal order}: For financial and time-series data, strictly employ chronological splits and forward-looking imputation to prevent look-ahead bias.
\end{enumerate}
\end{takeawaybox}
